\scanpage{37}
\[
H^{\unsure{1|n}}\bigl(X,\underbrace{\underline{\operatorname{Ext}}^{\unsure{1|n}}(f^*(B),G_{mX}})\bigr)
\]
\margin{\begin{tikzpicture}[x=1mm,y=1mm] \fill (0,4) ellipse (.12 and .17); \draw[line width=.3pt] (-.2,0) .. controls (.5,.05) and (.8,-.5) .. (.8,-1.2); \end{tikzpicture}}
\diagram{\begin{tikzcd}[ampersand replacement=\&,column sep=small,row sep=large,execute at end picture={
\path (\tikzcdmatrixname-2-1.north) -- (\tikzcdmatrixname-1-1.south)
 coordinate[pos=.15] (doublebase) coordinate[pos=.9] (doubletop);
\draw ([xshift=-.07cm]doublebase) .. controls ([xshift=-.12cm,yshift=.18cm]doublebase) and ([xshift=-.08cm,yshift=-.15cm]doubletop) .. ([xshift=-.07cm]doubletop);
\draw ([xshift=.04cm,yshift=.1cm]doublebase) .. controls ([xshift=.16cm,yshift=.09cm]doublebase) and ([xshift=.07cm,yshift=.17cm]doublebase) .. ([xshift=.07cm,yshift=.25cm]doublebase) -- ([xshift=.07cm,yshift=-.01cm]doubletop);
\draw ([xshift=-.10cm,yshift=-.13cm]doubletop) -- ([xshift=.05cm,yshift=.035cm]doubletop) -- ([xshift=.26cm,yshift=-.09cm]doubletop);
}]
\operatorname{Ext}_X(\overset{A}{\overset{\begin{tikzpicture}[baseline]\draw (0,0) .. controls (.25,.08) and (.55,.04) .. (.8,-.14);\end{tikzpicture}}{\unsure{f^*(B)|f_!(B)}}},G_{mX})
 \& {}\simeq\operatorname{Ext}_{\unsure{S|s}}(B,f_!(G_{mX}))\\
H^p(X,\underline{\operatorname{Ext}}^q_X(A,G_{mX})) \&
\end{tikzcd}}

\figure{\begin{tikzpicture}\draw (0,.12) .. controls (3,.25) and (6,.22) .. (10,0);\end{tikzpicture}}

\noindent\llap{\raisebox{-5.7cm}[0pt][0pt]{\begin{tikzpicture}[x=1cm,y=1cm]
\draw (0,5.65) .. controls (-.12,5.8) and (-.08,5.25) .. (-.06,4.9)
 .. controls (-.02,4.6) and (-.03,4.25) .. (-.17,4.17)
 .. controls (-.45,4.38) and (-.24,3.98) .. (-.13,4.08)
 .. controls (.02,3.85) and (.12,2.5) .. (.18,1.85)
 .. controls (.23,1.1) and (.24,.3) .. (.7,0);
\draw (.14,2.13) .. controls (.18,1.94) and (.2,1.84) .. (.31,1.79);
\draw (.27,.84)--(.52,.58);
\end{tikzpicture}}\hspace{.7em}}%
\begin{minipage}{\linewidth}
\unsure{Suites|Suite} \unsure{exactes|exacte} de Cohomologie. Critère de \unsure{rationalité};

\medskip
\unsure{Fibrés universels|Fibré universel}

\medskip
Dualité de \unsure{Cartier} [\unsure{voir} \unsure{Pic|Ric} \textendash\ \textendash\ \textendash]

\medskip
Th. \unsure{projectivité} \unsure{des|du} V.A. (Th. de \unsure{rigidité})

\medskip
Propriétés \unsure{fonctorielles|finicielles} des \unsure{groupes} \textendash\ \textendash\ \textemdash

\medskip
Propriétés différentielles. des \unsure{gpes|groupes} et fibrés

\medskip
Cas \unsure{particuliers}: \begingroup\setbox0=\hbox{fibrés linéaires,}\leavevmode
\rlap{\raisebox{-.45ex}{\begin{tikzpicture}[x=\wd0,y=1ex,baseline=0pt]
\draw (.025,0) .. controls (.35,.03) and (.8,.13) .. (1.09,.18);
\draw (.25,-.55) .. controls (.45,-.63) and (.75,-.35) .. (.97,-.38);
\end{tikzpicture}}}\box0\endgroup symplectiques, orthogonaux etc \textendash\ \textemdash\ \textendash
\end{minipage}

\medskip
\unsure{Sans} doute \unsure{par|pour} une \unsure{représentation|organisation} \unsure{lin.} de \unsure{gpes|groupes} ??

\figure{\begin{tikzpicture}\draw (0,0) .. controls (2,-.05) and (3.5,.14) .. (5,.07) .. controls (7.5,.02) and (8.5,.17) .. (10,.1);\end{tikzpicture}}
\[
\deleted{\underline{\operatorname{Hom}}(G_a,G_m)}
\]
\[
\mathbb Z.\quad\mathrel{\begin{tikzpicture}[x=1ex,y=1ex,baseline=.3ex]
\draw (.38,1.45) .. controls (.21,1.08) and (.03,.28) .. (.04,.02)
 .. controls (.08,-.18) and (.19,.22) .. (.34,.34)
 .. controls (.55,.21) and (.95,.3) .. (1.15,.34)
 .. controls (1.36,.4) and (1.56,.31) .. (1.66,.18);
\draw[line width=.6pt] (.35,.35) .. controls (.76,.2) and (1.15,.43) .. (1.53,.27);
\end{tikzpicture}}\quad\mathbb Z
\]
\margin{36}

